% ============================================================================
\section{Formulation du problème d'ordonnancement}
\label{sec:formulation}
% ============================================================================

\subsection{Données et variables de décision}

\begin{itemize}
  \item $T$ : l'ensemble des tâches du lot ; chaque tâche $i$ a un workload et une taille.
  \item $M$ : l'ensemble des VM ; chaque VM $j$ a un type, $k_j$ workers (= vCPU),
        une RAM et un prix horaire $\pi_j$.
  \item $x_{ij} \in \{0,1\}$ : vaut 1 si la tâche $i$ est exécutée sur la VM $j$.
  \item L'\textbf{ordre} des tâches dans la file de chaque VM, car l'orchestrateur
        les exécute dans l'ordre du fichier. L'ordre fait donc partie de la décision.
\end{itemize}

Une VM qui ne reçoit aucune tâche n'est pas allumée et ne coûte rien : choisir
les $x_{ij}$, c'est aussi choisir \textbf{quelles VM allumer}.

\subsection{Modèle de durée}

La durée prédite de la tâche $i$ sur la VM $j$ combine les trois mesures de la
section~\ref{sec:mesures} :
\begin{equation}
  p_{ij} = \underbrace{\tilde{t}_{ij}}_{\text{médiane du benchmark}}
  \times \underbrace{\gamma_j}_{\text{contention}}
  + \underbrace{\omega}_{\text{surcoût SSH}}
  \label{eq:pij}
\end{equation}
avec $\omega = 0{,}1$~s, $\gamma_j = 1{,}40$ pour les VM à 2~vCPU et $1{,}32$ pour
highcpu.

$\gamma_j$ est mesuré directement dans l'expérience workers
(section~\ref{sec:workers}) : c'est le temps de calcul de chaque tâche, relevé par
le workload lui-même quand tous les workers tournent, divisé par sa médiane de
benchmark (moyenne sur 24 tâches). Il ne se déduit pas du gain de débit
($4/2{,}40 \approx 1{,}67$ pour highcpu), qui inclut aussi le surcoût SSH (non
ralenti par la contention) et les workers inactifs en fin de lot. La valeur est
confirmée en réel : highcpu finit à $+2$ à $+5\,\%$ de sa prédiction dans toutes
nos exécutions (tableau~\ref{tab:precision}).

\subsection{Simulation, makespan et coût}

Pour évaluer une allocation sans l'exécuter, nous simulons exactement ce que fera
l'orchestrateur : chaque VM $j$ dispose de $k_j$ workers ; chaque tâche, dans
l'ordre de la file, est confiée au worker qui se libère le premier. On obtient
l'heure de fin $T_j$ de chaque VM, puis :
\begin{align}
  D &= \max_{j \in M} T_j \\
  C &= \sum_{j \,:\, \text{VM utilisée}} \max(60\ \text{s},\, T_j) \cdot \frac{\pi_j}{3600}
\end{align}
La facturation minimale de 60~s a une conséquence importante : \textbf{pour un
lot de moins d'une minute, le coût ne dépend que des VM allumées}, pas de la
façon dont on répartit les tâches. Au-delà, il dépend aussi du temps passé sur
chaque VM.

\subsection{Contraintes et objectif}

\begin{align}
  & \textstyle\sum_{j \in M} x_{ij} = 1 \quad \forall i \in T
    && \text{(chaque tâche exactement une fois)} \\
  & D \le 1{,}25 \cdot \Drr
    && \text{(délai, éq.~\ref{eq:deadline})} \\
  & k_j \cdot \max_{i : x_{ij}=1} \mathrm{RSS}_{ij} \le \mathrm{RAM}_j \quad \forall j
    && \text{(mémoire)} \\
  & \text{minimiser } S = 0{,}6\, C/\Crr + 0{,}4\, D/\Drr
    && \text{(éq.~\ref{eq:score})}
\end{align}
$\mathrm{RSS}_{ij}$ est la mémoire maximale mesurée de la tâche. Avec nos tailles,
la contrainte mémoire est toujours respectée (4 $\times$ 515~Mo $<$ 4~Go), mais
elle est vérifiée par le programme.

Il faut choisir une VM pour chaque tâche : avec 4 VM et 2000 tâches, cela fait
$4^{2000}$ possibilités (plus de $10^{1200}$), bien trop pour toutes les essayer.
Nous utilisons donc des \textbf{heuristiques} : des méthodes rapides qui
trouvent une bonne solution sans garantir la meilleure. Sur un lot de 9 tâches
($4^9 = 262\,144$ possibilités), on peut en revanche calculer l'optimum exact, ce
qui permet de mesurer l'écart de notre heuristique (section~\ref{sec:simu}).

\subsection{Validation du modèle}

Avant d'utiliser ce modèle pour comparer des algorithmes, nous avons vérifié
qu'il prédit bien la réalité. Pour cela, nous avons exécuté sur GCP un lot de
9 tâches réparti par Round Robin, puis comparé le makespan mesuré à celui prédit
par le modèle.

\begin{table}[H]
\centering\small
\begin{tabular}{lrrr}
\toprule
\textbf{Modèle} & \textbf{Makespan prédit} & \textbf{Makespan mesuré} & \textbf{Écart} \\
\midrule
Médiane du benchmark seule                             & 0,80 s & 1,32 s & réel $\times$ 1,65 \\
Avec contention et surcoût SSH (éq.~\ref{eq:pij})      & 1,22 s & 1,32 s & réel $\times$ 1,08 \\
\bottomrule
\end{tabular}
\caption{Makespan prédit et mesuré sur un lot de 9 tâches (\mesure{}).}
\label{tab:validation}
\end{table}

Avec la médiane seule, le modèle sous-estime le temps réel de 65\,\%. En
ajoutant la contention et le surcoût SSH, l'écart tombe à 8\,\%.

Ce lot n'a pas servi à régler les paramètres $\gamma$ et $\omega$, qui viennent
de l'expérience « workers » : il constitue donc un vrai test. Cette vérification
est intégrée aux tests automatiques du projet
(\code{TestPredictionMatchesRealRun}). La précision du modèle sur les lots de
500 et 2000 tâches est analysée en section~\ref{sec:reels}.
